\documentclass[source/paper/main.tex]{subfiles}
\begin{document}

\section{Model}
\subsection{Baseline model}

Time is discrete and indexed by $t$.
Each time period $t$, the project $i$ faces a latent number of problems $N_{i,t} \sim~ F_{it}(\cdot)$ where $F_{it}$ is some discrete probability mass function. 
Each problem is characterized by states. 
Within period $t$, let $X_{nt}^{(s)} \in \{\ell,o,r,m\}$ denote the state of problem $n$ after stage $s$ of the pull-request process, where $s \in \{0,1,2,3\}$ indexes stages within the period.

Each problem passes through the following stages:
\begin{enumerate}
    \item leads to an opened pull request,
    \item pull request is reviewed,
    \item pull request is merged.
\end{enumerate}

The state space is
\begin{itemize}
    \item $\ell$: latent problem for which no pull request has been opened,
    \item $o$: pull request opened to address the problem,
    \item $r$: opened pull request reviewed,
    \item $m$: reviewed pull request merged.
\end{itemize}
State $m$ is absorbing.

I will characterize the within-period stage transition structure for any problem in project $i$ at time $t$ by 
$$
\pi_{i,t} =
\begin{pmatrix}
1-\pi_{i,t}^o & \pi_{i,t}^o & 0 & 0 \\
0 & 1-\pi_{i,t}^r & \pi_{i,t}^r & 0 \\
0 & 0 & 1-\pi_{i,t}^m & \pi_{i,t}^m \\
0 & 0 & 0 & 1
\end{pmatrix}. 
$$
This characterization invokes Assumptions~\ref{asm:exchangeability} (Problem exchangeability), \ref{asm:independence} (Problem independence), and \ref{asm:memoryless} (Memoryless stage transitions).

Here, $\pi_{i,t}^o$ is the probability that a latent problem leads to an opened pull request, $\pi_{i,t}^r$ is the probability that an opened pull request is reviewed, and $\pi_{i,t}^m$ is the probability that a reviewed pull request is merged. 
The complementary probabilities $1-\pi_{i,t}^o$, $1-\pi_{i,t}^r$, and $1-\pi_{i,t}^m$ should be interpreted broadly. 
For example, a reviewed pull request could be left unmerged because the requested improvements were never implemented or because the reviewer deemed that the pull request's changes would not improve the software.

Hence, in each period $t$, there will be in expectation
\begin{enumerate}
    \item $N_{i,t}^o = N_{i,t} \cdot \pi_{i,t}^o$ opened pull requests,
    \item $N_{i,t}^r = N_{i,t}^o \cdot \pi_{i,t}^r$ reviewed pull requests, and
    \item $N_{i,t}^m = N_{i,t}^r \cdot \pi_{i,t}^m$ merged pull requests.
\end{enumerate}

My goal is to predict how organizational activity changes when a member leaves. To do that, I need a model of how individual actions contribute to organizational outcomes. 

\subsubsection{Opening, review, and merge as person-specific Bernoulli actions}
Let $\mathcal{J} = \{1, \ldots, J\}$ be the set of members. Each member $j \in \mathcal{J}$ contributes to organizational activity by opening, reviewing and merging pull requests.
For a given latent problem, at most one member can open the pull request. 
Any number of members can review an opened pull request.
For any reviewed pull request, at most one member can merge it.

For any latent problem in time period $t$, let $p_{j,t}^o$ denote the probability that member $j$ is the unique member who opens a pull request to address the problem. Hence, I can define the probability that no member opens a pull request,
$$ p_{0}^o = 1-\sum_{j\in\mathcal J} p_{j,t}^o. $$
which satisfies $ p_{j,t}^o \geq 0 $ and $ \sum_{j\in\mathcal J}p_{j,t}^o \leq 1. $

For each latent problem, the identity of the opening member is a draw from a categorical distribution over $\mathcal{J} \cup \{0\}$: member $j$ opens the pull request with probability $p_{j,t}^o$, and no member opens it with probability $p_{0}^o$. 
By assuming that draws are independent across problems, the joint count vector follows a multinomial distribution:
$$ \left( \left\{A_{j,t}^o \right\}_{j\in\mathcal J}, A_{0, t}^o \right) \mid N_{i,t} \sim \text{Multinomial}\left(N_{i,t}; \left \{p_{j,t}^o \right \}_{j\in\mathcal J}, p_{0}^o \right). $$

This additionally invokes Assumption~\ref{asm:problem-invariant} (Problem-invariant member probabilities).
Hence, the probability that a latent problem leads to an opened pull request is
$$ \pi_{i,t}^o = \sum_{j\in\mathcal J}p_{j,t}^o. $$


I model the quantity of pull requests reviewed by member $j$ in period $t$, $A_{j,t}^r$, as $N_{i,t}^o$ draws from a Binomial distribution:
$$A_{j,t}^r \mid N_{i,t}^o \sim \text{Binomial}(N_{i,t}^o, p_{j,t}^r),$$
where $p_{j,t}^r$ is member $j$'s probability of reviewing a given opened pull request. 
This also invokes Assumptions~\ref{asm:problem-invariant} (Problem-invariant member probabilities) and~\ref{asm:review-independence} (Independent review decisions).
Hence, the probability that an opened pull request is reviewed by at least one member is
$$ \pi_{i,t}^r = 1-\prod_{j\in\mathcal J}(1-p_{j,t}^r). $$


For any reviewed pull request at time $t$, let $p_{j,t}^m$ denote the probability that member $j$ is the unique member who merges it. Hence, I can define the probability that no member merges a pull request,
$$ p_{0}^m = 1-\sum_{j\in\mathcal J} p_{j,t}^m. $$
which satisfies $ p_{j,t}^m \geq 0 $ and $ \sum_{j\in\mathcal J}p_{j,t}^m \leq 1. $
For each reviewed pull request, the identity of the merging member is a draw from a categorical distribution over $\mathcal{J} \cup \{0\}$: member $j$ merges it with probability $p_{j,t}^m$, and no member merges it with probability $p_{0,t}^m$. By the same independence argument as for opening, the joint count vector follows a multinomial distribution:
$$ \left( \left \{A_{j,t}^m\right \}_{j\in\mathcal J}, A_{0, t}^m \right) \mid N_{i,t}^r \sim \text{Multinomial}\left(N_{i,t}^r; \left \{p_{j,t}^m\right \}_{j\in\mathcal J}, p_{0,t}^m \right). $$
Hence, the probability that a reviewed pull request is merged is $$ \pi_{i,t}^m = \sum_{j\in\mathcal J}p_{j,t}^m. $$
This also invokes Assumption~\ref{asm:problem-invariant} (Problem-invariant member probabilities).


\subsection{Prediction}
My goal is to predict, after a member $j^*$ departs in time period $t$, within an organization, the next-period $(t+1)$ quantity of
\begin{enumerate}
    \item pull requests opened ($N_{i,t+1}^o$),
    \item pull requests reviewed ($N_{i,t+1}^r$), and
    \item pull requests merged ($N_{i,t+1}^m$)
\end{enumerate}
I do that by fitting a model of the latent problem flow and member-level probabilities on organizational pre-period data, and then using the fitted model to predict counterfactual next-period activity after removing member $j^*$.

I impose two modeling choices to identify and estimate the latent problem flow:
\begin{enumerate}
    \item \textbf{Normalization of unobserved fraction}: $p_{0,t}^o = 0$ during the pre-periods, so every latent problem generates an opened pull request. This directly identifies the latent problem count from observed data, since $N_{i,t} = N_{i,t}^o$.
    \item \textbf{Time-invariant latent flow}: $F_{i,t} = F_i$, so pre-period opening counts $\{N_{i,t}^o\}_{t<0}$ are i.i.d.\ draws from $F_i$, which I estimate by maximum likelihood; I denote the fitted distribution $\hat{F}_i$ and its mean $\hat{N}_i$.
\end{enumerate}

I evaluate two fitting approaches, each applied organization-by-organization on pre-period opening counts by maximum likelihood:
\begin{enumerate}
    \item \textbf{Poisson/ZIP}: fit both Poisson and ZIP for each organization; select between them by comparing the sum of log likelihoods across periods.
    \item \textbf{Negative binomial}: fit a negative binomial for each organization.
\end{enumerate}
To calculate the log likelihood for each fitting approach, for each time period, I fit the data using the other four time periods, and then evaluate the log likelihood of the fitted distribution on the held-out period's data; I sum across the five held-out periods to obtain a total log likelihood for each distribution.
Since pre-period lengths are the same across organizations, and each organization uses the same data (just a different distribution) to fit the latent problem flow, the log likelihood of each fitted distribution is comparable within organization. 
To compare approaches, I will report the cross-organization distribution of pairwise log likelihood differences (e.g., approach 1 minus approach 2 per organization), with the median log likelihood difference as a summary statistic. 
I will also plot the distribution of log likehoods across organizations for each fitting approach. 
I also store the empirical versus fitted PMF for each organization as an auxiliary visual check.

I impose two additional modeling choices governing member probabilities:
\begin{enumerate}
    \item \textbf{Stationary member probabilities}: $p_{j,t}^s = p_j^s$ for all $t$ and all $s$. Under stationarity, pre-period observations pool cleanly across periods for estimation, and the counterfactual after departure is the model evaluated at the same probabilities with member $j^*$ removed.
    \item \textbf{Fixed membership}: the membership set $\mathcal{J}$ is fixed across periods, with no entry or exit other than the focal departure $j^*$.
\end{enumerate}
Note that the remaining members' probabilities are unchanged after the departure. This is follows directly from the Assumption~\ref{asm:problem-invariant} (Problem-invariant member probabilities).

I use one estimation approach for member probabilities:
\begin{enumerate}
    \item \textbf{Pooled pre-period shares}: all member probabilities are identified by pre-period shares, requiring no numerical optimization. For opening and merging, the MLE under a stationary multinomial is $\hat{p}_j^o = \frac{\sum_t A_{j,t}^o}{\sum_t N_{i,t}^o}$ and $\hat{p}_j^m = \frac{\sum_t A_{j,t}^m}{\sum_t N_{i,t}^r}$. For reviewing, the Binomial MLE is $\hat{p}_j^r = \frac{\sum_t A_{j,t}^r}{\sum_t N_{i,t}^o}$.
\end{enumerate} 
For reviewing, I can also check whether the implied aggregate review probability $\hat{\pi}^r = 1-\prod_{j\in\mathcal{J}}(1-\hat{p}_j^r)$ matches the observed per-period review rate $\tilde{\pi}_t^r = N_{i,t}^r / N_{i,t}^o$, plotting the cross-organization distribution of the mean squared deviation between $\hat{\pi}^r$ and $\tilde{\pi}_t^r$ across pre-periods, and reporting the median and various percentiles as a summary statistic.


%MODELLING choices I will not make now but might be good to incorporate
%- model entry/exit from the org 
%- Is there a way to penalize big p's?
%- find ways to loosen assumptions - I feel like I have 5 structural ones + the modelling ones above.  - ACTUALLY UPON FURTHER REFLECTION, I THINK I SHOULD THINK ABOUT HOW I"M GOING TO VARY THE MODEL BASED ON WHAT I"M FAILING TO PREDICT WELL. THE TRICKY THING HERE IS THAT I ONLY HAVE ONE OBSERVATION FOR "HOW WELL I"M PREDICTING". SO I HAVE TO THINK ABOUT HOW DISTRIBUTIONALLY ACROSS ORGS MODELLING CHOICES IMPROVE PREDICTIVE PERFORMANCE??? 


Using the fitted quantities $\hat{N}_i$, $\hat{p}_j^s$, I define the counterfactual stage-$s$ completion probability after removing member $j^*$ as $\hat{\pi}_i^{s,(-j^*)}$, and the predicted next-period count of stage-$s$ activity as $\hat{N}_{i,t+1}^s$.  
The counterfactual probability of a pull request being opened from a latent problem is  
$$ \hat{\pi}_i^{o,(-j^*)} = \sum_{j\in\mathcal J\setminus\{j^*\}}\hat{p}_j^o, $$
so the expected counterfactual number of pull requests opened is 
$$ \hat{N}_{i,t+1}^o = \hat{N}_i\,\hat{\pi}_i^{o,(-j^*)}. $$
For the initial prediction, I am hence assuming that the probabilities of all members $j$ who showed up only in periods following the departure are 0. The counterfactual review probability is then
$$ \hat{\pi}_i^{r,(-j^*)} = 1-\prod_{j\in\mathcal J\setminus\{j^*\}}(1-\hat{p}_j^r), $$ 
so the expected counterfactual number of pull requests reviewed is 
$$ \hat{N}_{i,t+1}^r = \hat{N}_i\,\hat{\pi}_i^{o,(-j^*)}\hat{\pi}_i^{r,(-j^*)}. $$
The counterfactual merge probability is 
$$ \hat{\pi}_i^{m,(-j^*)} = \sum_{j\in\mathcal J\setminus\{j^*\}}\hat{p}_j^m, $$ 
so the expected counterfactual number of pull requests merged is 
$$ \hat{N}_{i,t+1}^m = \hat{N}_i\,\hat{\pi}_i^{o,(-j^*)}\hat{\pi}_i^{r,(-j^*)}\hat{\pi}_i^{m,(-j^*)}. $$


\subsubsection{Evaluating predictive performance}

I will evaluate predictive performance by comparing the predicted next-period activity to the observed next-period activity for each stage $s$. 
For the quantity of pull requests opened, I will compare $\hat{N}_{i,t+1}^o$ to $N_{i,t+1}^o$, plot the cross-organization distribution of prediction errors, in percentage points relative to the prediction and report summary statistics. 

Since prediction errors from the quantity of pull requests opened will propagate to the quantity of pull requests reviewed and merged, to compare the quantity of pull requests reviewed and merged, I will instead compare the predicted stage completion probabilities $\hat{\pi}_i^{r,(-j^*)}$ and $\hat{\pi}_i^{m,(-j^*)}$ to the observed stage completion probabilities $\tilde{\pi}_{i,t+1}^r = N_{i,t+1}^r / N_{i,t+1}^o$ and $\tilde{\pi}_{i,t+1}^m = N_{i,t+1}^m / N_{i,t+1}^r$, plotting the cross-organization distribution of prediction errors in raw percentage points and reporting summary statistics.


\iffalse



\subsubsection{Putting everything together}


\label{sec:stage_decomposition_baseline}

The stage-completion probabilities defined above together imply
$$
E[Q_t^m \mid N_{i,t}]
=
N_{i,t} \pi_{i,t}^o \pi_{i,t}^r \pi_{i,t}^m.
$$

Now define member $j$'s marginal contribution at each stage as the change in the stage-completion probability when $j$ is removed:
$$
\Delta_o^{(j)}
\equiv
\pi_{i,t}^o - \pi_{i,t}^{o,(-j)}
=
p_{j,t}^{o}\prod_{j\in\mathcal J\setminus\{j\}}(1-p_{j,t}^{o}),
$$
$$
\Delta_r^{(j)}
\equiv
\pi_{i,t}^r - \pi_{i,t}^{r,(-j)}
=
p_{j,t}^{r}\prod_{j\in\mathcal J\setminus\{j\}}(1-p_{j,t}^{r}),
$$
$$
\Delta_m^{(j)}
\equiv
\pi_{i,t}^m - \pi_{i,t}^{m,(-j)}
=
p_{j,t}^{m}\prod_{j\in\mathcal J\setminus\{j\}}(1-p_{j,t}^{m}).
$$

Suppose member $j$ leaves before period $t+1$ and is not replaced, and suppose $N_{i,t}=N_{t+1}=N$. Then
$$
E[Q_{t+1}^m \mid N]
=
N \pi_{i,t}^{o,(-j)} \pi_{i,t}^{r,(-j)} \pi_{i,t}^{m,(-j)}
=
N(\pi_{i,t}^o-\Delta_o^{(j)})(\pi_{i,t}^r-\Delta_r^{(j)})(\pi_{i,t}^m-\Delta_m^{(j)}).
$$
Therefore,
$$
E[Q_{t+1}^m \mid N] - E[Q_t^m \mid N]
=
-N\Big[
\Delta_o^{(j)}(\pi_{i,t}^r-\Delta_r^{(j)})(\pi_{i,t}^m-\Delta_m^{(j)})
+\Delta_r^{(j)}\pi_{i,t}^o(\pi_{i,t}^m-\Delta_m^{(j)})
+\Delta_m^{(j)}\pi_{i,t}^o \pi_{i,t}^r
\Big].
$$

The three terms correspond to member $j$'s marginal contribution at the opening, review, and merge stages:
\begin{enumerate}
    \item $-N\,\Delta_o^{(j)}(\pi_{i,t}^r-\Delta_r^{(j)})(\pi_{i,t}^m-\Delta_m^{(j)})$ captures the loss from problems that would have been opened because member $j$ acts at the opening stage, and that would then have been reviewed and merged without member $j$'s involvement at those later stages.

    \item $-N\,\pi_{i,t}^o\Delta_r^{(j)}(\pi_{i,t}^m-\Delta_m^{(j)})$ captures the loss from pull requests that would have been opened (by anyone), then reviewed because member $j$ acts at the review stage, and then merged without member $j$'s involvement at the merge stage.

    \item $-N\,\pi_{i,t}^o \pi_{i,t}^r \Delta_m^{(j)}$ captures the loss from pull requests that would have been opened and reviewed (by anyone), and then merged because member $j$ acts at the merge stage.
\end{enumerate}
\fi 

\subsection{Structural assumptions}

\begin{assumption}[Problem exchangeability]
\label{asm:exchangeability}
Problems are indistinguishable: for any two problems $n$ and $n'$ and any states $x,y \in \{\ell,o,r,m\}$,
$$
\Pr(X_{n,t}^{(s+1)}=y \mid X_{n,t}^{(s)}=x)
=
\Pr(X_{n',t}^{(s+1)}=y \mid X_{n',t}^{(s)}=x).
$$
Conditional on being in the same state, any two problems have the same probability of advancing to the next stage.
\end{assumption}

\begin{assumption}[Problem independence]
\label{asm:independence}
Conditional on current states, problems evolve independently of one another:
$$
\Pr(X_{1,t}^{(s+1)}=x_1,\dots,X_{N_{i,t},t}^{(s+1)}=x_{N_{i,t}}\mid X_{1,t}^{(s)},\dots,X_{N_{i,t},t}^{(s)})
=
\prod_{n=1}^{N_{i,t}}\Pr(X_{n,t}^{(s+1)}=x_n\mid X_{n,t}^{(s)}).
$$
\end{assumption}

\begin{assumption}[Memoryless stage transitions]
\label{asm:memoryless}
A problem's next stage depends only on its current stage, not its history:
$$
\Pr(X_{n,t}^{(s+1)}=y \mid X_{n,t}^{(s)},X_{n,t}^{(s-1)},\dots)
=
\Pr(X_{n,t}^{(s+1)}=y \mid X_{n,t}^{(s)}).
$$
\end{assumption}

\begin{assumption}[Problem-invariant member probabilities]
\label{asm:problem-invariant}
A member's probability of acting does not vary across problems within a period --- $p_{j,t}^s$ is the same regardless of which particular problem is being acted on. This is distinct from Assumption~\ref{asm:exchangeability}, which describes the distribution problems are drawn from: even if problems differed objectively, this rules out members selectively targeting particular problems.
\end{assumption}

Assumptions \ref{asm:exchangeability}--\ref{asm:problem-invariant} imply that every problem follows the same stage-transition matrix $P$ so that, for every problem $n$,
$$
\Pr(X_{n,t}^{(s+1)}=y \mid X_{n,t}^{(s)}=x)=P_{xy}.
$$

\begin{assumption}[Independent review decisions]
\label{asm:review-independence}
Conditional on an opened pull request, each member $j$ independently decides whether to review it. Whether member $j$ reviews a pull request does not affect whether member $j'$ does so. Since each member's decision whether to review a pull is a Bernoulli trial, independence across members implies that the quantity of pull requests reviewed by member $j$ follows a Binomial distribution. 
\end{assumption}

\iffalse
\subsection{Extensions}

\begin{enumerate}
    \item \textbf{Relax Assumption 1.} Problems need not come from an identical distribution. For example, the identity of the pull request opener may affect the probability that the pull request is eventually merged.

    \item \textbf{Relax Assumption 2.} Problems need not evolve independently. For example, one pull request being merged may increase the likelihood that followup problems are opened or closed.

    \item \textbf{Relax Assumption 4.} Member-specific probabilities may vary over time. For example, one could model $\pi_{i,t}^{o}$, $\pi_{i,t}^{r}$, and $\pi_{i,t}^{m}$ as draws from member-specific distributions, although this may require more data than are currently available.

    \item \textbf{Relax Assumption 5.} Members' actions within a stage need not be independent. For example, reviewers may coordinate so that one person's review makes another person's review less likely, or a member's departure may induce other members to increase their probabilities of acting. This extension would replace the independent-Bernoulli benchmark with a model that allows for coordination and/or organizational adaptation.
\end{enumerate}
\fi
\end{document}
